\chapter{Semantic Geometry}\label{ch:geometry}

\begin{workplan}{\Keep}
\wpgoal{Keep the geometry and make its dispatch rule --- which metric for
which job --- the chapter's organising principle rather than a remark.}

\wpsource{This chapter's existing text;
\texttt{MATH\_TOOLBOX.md}~IV.1 (the estimation/comparison dispatch), IV.2
(Cone--Bures: kept as a metric, its closeness claim refuted), IV.3 (what a
concept is), I.6 (the variance floor --- the largest measured downstream
effect in the V0 record).}

\wpsteps{
  \item \textbf{Lead with the dispatch rule.} \textsc{estimation} uses
    information/Fisher--Rao geometry, where a theorem exists (precision-
    weighted fusion is a genuine MLE); \textsc{comparison} uses
    Bures--Wasserstein transport geometry, where a modelling choice exists.
    They genuinely disagree --- by Agueh--Carlier the Wasserstein barycenter
    of Gaussians has the \emph{plain} weighted mean, not the precision-weighted
    one. Never blend comparison into estimation.
  \item \textbf{Own the seam.} $\pi_v$ fuses in information geometry but its
    output is immediately compared in transport geometry by $\omega$; $\pi_v$
    is \emph{not} the transport barycenter of its own fine complex. This is a
    paragraph the book owes a reader, not a bug. Write it.
  \item \textbf{Record the rejected repair.} Training a blend weight between
    the two geometries is \Refuted, on four grounds worth keeping because they
    generalise: it replaces a choice-with-a-reason by a fitted-constant-with-
    none; no training signal exists; one scalar wants a sweep, not an
    optimizer; and the two uses are asymmetric, so blending damages a proven
    MLE for a fitted number.
  \item \textbf{Upgrade what a concept \emph{is}.} A Gaussian stalk's level
    sets are ellipsoids, hence convex, so the SPD stalk already \emph{is} a
    convex conceptual region in Gärdenfors' sense --- concepts have extent, and
    treating one as a point discards exactly the structure this geometry
    builds. State this as the chapter's conclusion.
  \item \textbf{Keep the variance-floor fix and its numbers.} The correct
    invariant is ``floor carries trace $O(1)$'' and not ``shared epsilon'';
    getting this wrong once cost a $47\times$ slowdown, and fixing it moved the
    epistemic term from $229.76$ to $0$ and made the transport term non-inert.
    \Measured. It is the best worked example in the book of a hardcoded
    constant silently flattening a whole geometry.
  \item \textbf{Keep Cone--Bures with its refutation attached.} It survives as
    a genuine metric (proved, measured, bounded, and it deletes a hand-set
    threshold). Its ``tracks true Hellinger--Kantorovich to under 1\%'' claim
    does \emph{not} survive real text. Keep both facts adjacent; report the
    diagnostic ratio as telemetry and never as a claim.
}

\wpdone{A reader can state which metric to use for a given job and why, and
can explain the seam between them without treating it as an oversight.}
\end{workplan}

Part~I built the shape and Part~II measured disagreement across it. This
part supplies the geometry \emph{inside} an organ: what a concept is,
how many concepts collapse into the single vector a stalk holds, and how
far apart two uncertain concepts are.

\section{A concept is a Gaussian}

\begin{intu}\Intu
A concept is not a point. It is a point plus an admission of how
uncertain that point is, and --- more subtly --- of the \emph{directions}
in which it is uncertain. A concept may be sharply located along one
axis of meaning and vague along another. A Gaussian with a structured
covariance is the smallest object that records all three.
\end{intu}

\begin{definition}[Concept]\label{def:concept-gauss}
A \emph{concept} is a Gaussian $\mathcal N(\mu_i,\Sigma_i)$ on $\Rd$
with low-rank-plus-isotropic covariance
\begin{equation}\label{eq:cov}
  \Sigma_i=U_iU_i^\top+D_i I_d,
  \qquad U_i\in\R^{d\times k_i},\ D_i>0 .
\end{equation}
Here $\mu_i$ is the embedding, $U_iU_i^\top$ is anisotropic learned
spread, and $D_i$ is an isotropic noise floor --- the epistemic humility
of the concept.
\end{definition}

The rank $k_i$ is kept small (in practice $k_i\le8$). This is not a
performance compromise: it is what makes every computation below reduce
to a $k\times k$ problem rather than a $d\times d$ one, and it is the
reason the geometry is affordable at all.

\section{The coarse-graining $\pi_v$}

\begin{intu}\Intu
An organ holds many concepts but exposes one vector to the coarse
complex. Which vector? Not the plain average: a near-certain concept
should dominate a vague one. The right notion of weight is
\emph{precision} --- inverse variance --- and the resulting formula is
exactly how independent noisy measurements of a single quantity are
optimally combined.
\end{intu}

\begin{construction}[The coarse-graining map]\label{con:piv}
For an organ $v$ with active concepts
$\{\mathcal N(\mu_i,\Sigma_i)\}_{i=1}^n$, the coarse stalk value is
\begin{equation}\label{eq:piv}
  \boxed{\ \pi_v\;=\;
  \Bigl(\sum_{i=1}^n\Sigma_i^{-1}\Bigr)^{-1}
  \Bigl(\sum_{i=1}^n\Sigma_i^{-1}\mu_i\Bigr)\ }
\end{equation}
If $v$ has no active concept then $\pi_v=\unk$, and $v$ is excluded from
the computation of $\rhoscore$ rather than contributing a zero vector.
\end{construction}

\begin{proposition}[$\pi_v$ is the maximum-likelihood fusion]\label{prop:pivmle}
\Proved\ Suppose $\theta\in\Rd$ is a latent vector and each concept is an
independent observation $\mu_i=\theta+\varepsilon_i$ with
$\varepsilon_i\sim\mathcal N(0,\Sigma_i)$. Then the maximum-likelihood
estimate of $\theta$ is \eqref{eq:piv}.
\end{proposition}
\begin{proof}
The log-likelihood is
$\ell(\theta)=-\tfrac12\sum_i(\mu_i-\theta)^\top\Sigma_i^{-1}(\mu_i-\theta)
+\text{const}$.
Its gradient is
$\nabla_\theta\ell=\sum_i\Sigma_i^{-1}(\mu_i-\theta)$, and setting it to
zero gives $\bigl(\sum_i\Sigma_i^{-1}\bigr)\theta=\sum_i\Sigma_i^{-1}\mu_i$.
The matrix $\sum_i\Sigma_i^{-1}$ is positive definite (each
$\Sigma_i^{-1}$ is, since $D_i>0$), hence invertible, and $\ell$ is
strictly concave, so the stationary point is the unique maximum.
\end{proof}

\begin{remark}[Three properties worth naming]
The fusion \eqref{eq:piv} is \emph{order-independent} (both sums are
commutative), \emph{non-destructive} (every concept contributes to both
the precision sum and the mean; none is discarded), and \emph{$n$-ary}
rather than iterated-pairwise. These are not incidental: an
order-dependent or lossy summary would make the coarse measurement
depend on the sequence in which an organ happened to acquire its
concepts.
\end{remark}

\subsection*{Two computational paths, both exact}

\begin{itemize}
  \item \textbf{Isotropic fast path.} When all $U_i$ are empty,
  $\Sigma_i^{-1}=D_i^{-1}I$ and \eqref{eq:piv} collapses to the scalar
  precision-weighted average
  \begin{equation}\label{eq:isopath}
    \pi_v=\frac{\sum_i D_i^{-1}\mu_i}{\sum_i D_i^{-1}},
  \end{equation}
  computable in $O(nd)$ with no matrix inversion.

  \item \textbf{General path.} With $U_i\neq\varnothing$, invert each
  $\Sigma_i$ by the Sherman--Morrison--Woodbury identity
  \begin{equation}\label{eq:woodbury}
    \Sigma_i^{-1}=D_i^{-1}I
    -D_i^{-1}U_i\bigl(I_{k_i}+D_i^{-1}U_i^\top U_i\bigr)^{-1}U_i^\top D_i^{-1},
  \end{equation}
  accumulate into a dense information matrix, and solve once, at
  $O(ndk+d^3)$.
\end{itemize}

\begin{lemma}[Woodbury for this covariance]\label{lem:woodbury}
\Proved\ \eqref{eq:woodbury} holds for $\Sigma=UU^\top+DI$ with $D>0$.
\end{lemma}
\begin{proof}
The Woodbury identity states
$(A+UCV)^{-1}=A^{-1}-A^{-1}U(C^{-1}+VA^{-1}U)^{-1}VA^{-1}$. Take
$A=DI$, $C=I_k$, $V=U^\top$. Then $A^{-1}=D^{-1}I$ and
$C^{-1}+VA^{-1}U=I_k+D^{-1}U^\top U$, which is positive definite and
hence invertible. Substituting gives \eqref{eq:woodbury}. The cost is
dominated by the $k\times k$ inverse, not by anything of size $d$.
\end{proof}

\begin{remark}[The floor is not optional]
A lower bound $D_i\ge10^{-6}$ is imposed. Without it, a concept reported
as near-certain produces $1/D_i\approx10^{16}$ and the information
matrix loses all contribution from every other concept to rounding. The
floor pins $\pi_v$ strongly but finitely. Numerically the general path
was checked against a dense $\bigl(\sum_i\Sigma_i^{-1}\bigr)^{-1}$ solve
and agreed to $2.2\times10^{-16}$; equal precisions reduce to the plain
mean, and a $9{:}1$ precision ratio pulls $\pi_v$ to the confident
concept, as required.
\end{remark}

\begin{remark}[This map is the two-level glue]\label{rem:glue}
After execution, each organ's coarse stalk is \emph{overwritten} by its
own $\pi_v$. That is what makes the coarse $\rhoscore$ and the fine
geometry one measurement at two resolutions rather than two unrelated
numbers. It also means $\pi_v$ is lossy in a load-bearing way: whatever
the fine complex knew that does not survive \eqref{eq:piv} is invisible
to every coarse diagnostic.
\end{remark}

\section{Distance between uncertain concepts}

\begin{intu}\Intu
How far apart are two concepts? Comparing centres alone throws away the
spreads. The $2$-Wasserstein --- optimal-transport --- distance between
Gaussians accounts for both, and is a genuine metric rather than a
heuristic score.
\end{intu}

\begin{definition}[$2$-Wasserstein between Gaussians]\label{def:bures}
\begin{equation}\label{eq:bures}
  W_2^2\bigl(\mathcal N(\mu_1,\Sigma_1),\mathcal N(\mu_2,\Sigma_2)\bigr)
  =\norm{\mu_1-\mu_2}^2+\Tr\Sigma_1+\Tr\Sigma_2
   -2\Tr\Bigl[\bigl(\Sigma_1^{1/2}\Sigma_2\Sigma_1^{1/2}\bigr)^{1/2}\Bigr].
\end{equation}
The covariance part alone is the squared Bures--Wasserstein distance
$\BW^2$ on $\SPD{d}$.
\end{definition}

\begin{proposition}[Isotropic reduction to $O(k^3)$]\label{prop:buresiso}
\Proved\ If $\Sigma_1=D_1I$ is isotropic and
$\Sigma_2=U_2U_2^\top+D_2I$ with $U_2\in\R^{d\times k}$, then, writing
$\{\sigma_j^2\}_{j=1}^k$ for the eigenvalues of the $k\times k$ matrix
$U_2^\top U_2$,
\begin{equation}\label{eq:buresreduced}
  W_2^2=\norm{\mu_1-\mu_2}^2+\bigl(\Tr\Sigma_1+\Tr\Sigma_2\bigr)
  -2\sqrt{D_1}\Bigl[\sum_{j=1}^k\sqrt{\sigma_j^2+D_2}+(d-k)\sqrt{D_2}\Bigr].
\end{equation}
\end{proposition}
\begin{proof}
Since $\Sigma_1^{1/2}=\sqrt{D_1}I$, we get
$\Sigma_1^{1/2}\Sigma_2\Sigma_1^{1/2}=D_1\Sigma_2$ and hence
$\bigl(D_1\Sigma_2\bigr)^{1/2}=\sqrt{D_1}\,\Sigma_2^{1/2}$, so the cross
term is $2\sqrt{D_1}\Tr\Sigma_2^{1/2}$. The eigenvalues of
$\Sigma_2=U_2U_2^\top+D_2I$ are $\sigma_j^2+D_2$ on $\im U_2$ and $D_2$
with multiplicity $d-k$ on $\ker U_2^\top$, because the nonzero
eigenvalues of $U_2U_2^\top$ coincide with those of $U_2^\top U_2$.
Summing $\sqrt{\lambda}$ over the spectrum gives
\eqref{eq:buresreduced}. Only a $k\times k$ eigendecomposition is
needed.
\end{proof}

This is exact for isotropic $\Sigma_1$, not a diagonal approximation.

\section{The two terms are not commensurable}

\begin{construction}[Reporting the split]\label{con:w2split}
$W_2^2$ is a sum of two quantities measuring different things, and they
are returned separately rather than added:
\begin{equation}\label{eq:w2split}
  W_2^2=
  \underbrace{\norm{\mu_1-\mu_2}^2}_{\text{semantic: \emph{what} they say}}
  +
  \underbrace{\Tr\Sigma_1+\Tr\Sigma_2
   -2\Tr\bigl[(\Sigma_1^{1/2}\Sigma_2\Sigma_1^{1/2})^{1/2}\bigr]}
   _{\text{epistemic: \emph{how sure} they are}} .
\end{equation}
\end{construction}

\begin{proposition}[Degeneration in the isotropic regime]\label{prop:w2hazard}
\Proved\ If both concepts are isotropic, $\Sigma_i=D_iI$ with $U_i$
empty, then
\begin{equation}\label{eq:w2iso}
  W_2^2=\norm{\mu_1-\mu_2}^2+d\bigl(\sqrt{D_1}-\sqrt{D_2}\bigr)^2 .
\end{equation}
In particular $D_1=D_2$ implies $W_2^2=\norm{\mu_1-\mu_2}^2$ exactly.
\end{proposition}
\begin{proof}
Set $k=0$ in Proposition~\ref{prop:buresiso}. Then
$\Tr\Sigma_1+\Tr\Sigma_2=d(D_1+D_2)$ and the cross term is
$2\sqrt{D_1}\cdot d\sqrt{D_2}$, so the covariance part equals
$d\bigl(D_1+D_2-2\sqrt{D_1D_2}\bigr)=d(\sqrt{D_1}-\sqrt{D_2})^2$.
\end{proof}

\begin{hon}\Hon
Proposition~\ref{prop:w2hazard} recorded two uncomfortable facts. Both
were true when written; both were subsequently repaired. The diagnosis
is retained because it is what forced the repair and because the
structural lesson survives it.

\emph{(i) The machinery was inert.} Every grown concept carried the same
prior $D=1.0$ with empty $U$, so $D_1=D_2$ and $W_2^2$ was
\emph{exactly} squared Euclidean distance --- an optimal-transport
apparatus with a Wasserstein-shaped proof attached and nothing to do.
After the repair, $D$ is the shrinkage floor of
Remark~\ref{rem:stalkfloor} and varies with effective sample size and
with rank, so $D_1\neq D_2$ in general. On two stalks with equal means
and ranks $0$ and $1$, the epistemic term moves from $0$ to $0.206$: the
geometry is live.

\emph{(ii) It would have become actively harmful the moment calibration
was fixed.} The semantic term is $d$-\emph{intensive} --- bounded by $4$
for unit-normalised embeddings, whatever $d$ is --- while the epistemic
term is $d$-\emph{extensive}. At $d=384$, a prior $D_1=1.0$ against a
confident concept with $D_2=-\ln0.95$ gave an epistemic term of roughly
$230$ against a semantic term of $2$: about $115$ to $1$, driven
entirely by a self-reported confidence that the architecture elsewhere
declines to trust. Merge, split and consolidation would then have been
governed by hallucinated confidences. This is a latent defect that
switches on precisely when the system is \emph{improved}, which is the
worst kind.

The lesson survives both fixes and is why the terms are reported apart:
summing a $d$-intensive and a $d$-extensive quantity and thresholding
the sum is a dimensional error, not a modelling preference.
\end{hon}

\begin{exercise}
Verify the $115{:}1$ figure. Take $d=384$, $D_1=1.0$, $D_2=-\ln 0.95$,
both isotropic, and evaluate \eqref{eq:w2iso}. Then repeat with the
shrinkage floor $D=O(1/d)$ of Remark~\ref{rem:stalkfloor} and confirm
the epistemic term drops to order $1$. The point of doing this by hand
is that the failure is visible in the exponents, before any numbers are
substituted.
\end{exercise}

\section{Fusion and its entropy}

\begin{definition}[Precision-weighted fusion]\label{def:fusion}
The fusion of two concepts is the information-form product
\begin{equation}\label{eq:merge}
  \Lambda_i=\Sigma_i^{-1},\qquad
  \Sigma_\star=(\Lambda_1+\Lambda_2)^{-1},\qquad
  \mu_\star=\Sigma_\star(\Lambda_1\mu_1+\Lambda_2\mu_2),
\end{equation}
which is \eqref{eq:piv} for $n=2$. Full covariances are fused via
Woodbury \eqref{eq:woodbury}, never element-wise.
\end{definition}

\begin{proposition}[Fused entropy by the matrix determinant lemma]\label{prop:entropy}
\Proved\ The differential entropy
$h=\tfrac12\ln\bigl((2\pi e)^d\det\Sigma_\star\bigr)$ is computable from
a $k\times k$ determinant: for $\Sigma=UU^\top+DI$,
\begin{equation}\label{eq:detlemma}
  \det\Sigma=D^d\det\bigl(I_k+D^{-1}U^\top U\bigr).
\end{equation}
\end{proposition}
\begin{proof}
The matrix determinant lemma gives
$\det(A+UV^\top)=\det A\cdot\det(I+V^\top A^{-1}U)$. With $A=DI$ and
$V=U$ we have $\det A=D^d$ and $V^\top A^{-1}U=D^{-1}U^\top U$.
\end{proof}

\begin{remark}[Where confidence actually goes]\label{rem:stalkfloor}
\Implemented\ An earlier design set the noise floor from the generating
model's self-reported confidence via $D=-\ln c$, with a named prior
$D=1.0$ when $c$ was unavailable. Both were removed. The floor is now a
shared shrinkage floor
\[
  D=\varepsilon+\frac{\kappa}{d\,(n_{\mathrm{eff}}+\kappa)}=O(1/d),
\]
independent of any reported confidence. The scaling is the whole point:
$D$ feeds a $d$-extensive Bures term, so a floor of order $1$
contributes a trace of order $d$ and swamps a semantic term bounded by
$4$. At $d=384$ the two retired forms give isotropic traces of about
$19.7$ and $384$ respectively; the shrinkage floor gives $0.88$.

Confidence is \emph{relocated, not discarded}: its home is the edge
precision of Construction~\ref{con:weighted}, entering as
$s_e=-\ln c_e/d$ so that its trace contribution is $O(1)$. There it is
dimensionally harmless and load-bearing; as an isotropic stalk variance
it was neither. Against the deployed model, token log-probabilities are
unavailable, so $c$ is genuinely unknown and the engine warns once
rather than fabricating a value. This is why $\rhoscore$, and not
self-reported confidence, is the primary signal.
\end{remark}

\begin{hon}\Hon
The pattern in this chapter is worth extracting, because it recurs. Each
of the three defects --- inert geometry, dimensional mismatch,
fabricated confidence --- was invisible to every test that checked
whether the code \emph{ran}, and visible immediately to a check of
whether the quantities being combined had the same units. Dimensional
analysis is the cheapest correctness tool available here and it was
applied late.
\end{hon}
